\chapter{Symbol Glossary and Guide to the Literature}
\label{app:glossary}

\chapterorient{This closing appendix is a reference, not a narrative. Its first
half is a glossary: every symbol the book uses often, grouped by the world it
lives in---fields, operators, function spaces, numerical methods, the
pseudo-language, and the layered methodology---with a pointer to the chapter that
introduced it. Its second half is a guide to the literature, a topic-organized
map of the sources the book leans on, each with a one-line note on what to read
it for. The notation primer at the front of the book introduced these symbols a
few at a time; here they are collected so you can look one up. When a symbol is
unfamiliar, start here.}

This appendix expands the front-matter notation primer (and its
\cref{tab:symbols}) into a full reference. The mathematical symbols are set the
usual way; the pseudo-language symbols obey the one standing convention of the
book---they belong in a shaded box, never inside a formula, because the
shape-group sigil \texttt{\$} clashes with mathematics typesetting. In the
tables below, pseudo-language identifiers are set in \texttt{typewriter}; where a
shape carries the sigil we write it as \texttt{Tensor[(S:\,\dots)]} and discuss
the \texttt{\$S} back-reference form in prose.

\section{Symbol and notation glossary}

The glossary is split into six categories, one cluster of tables each: physical
fields and quantities; discrete operators and matrices; function spaces and the
finite-element vocabulary; the symbols of the numerical methods; the
pseudo-language and calculus; and the layered methodology. Each row gives the
symbol, its meaning, and a cross-reference to the chapter where it is introduced
and used.

\subsection{Fields and physical quantities}

These are the quantities Maxwell's equations relate, the materials they move
through, and the scalar readouts a simulation finally reports. Vector fields are
upright bold; material coefficients and scalars are italic.

\begin{table}[ht]
\centering
\small
\caption[Fields and material coefficients]{Electromagnetic fields, potentials,
and the constitutive coefficients of the medium.}
\label{tab:gl-fields}
\begin{tabular}{@{}lll@{}}
\toprule
Symbol & Meaning & Where introduced \\
\midrule
$\vE$ & electric field & \cref{ch:maxwell} \\
$\vB$ & magnetic flux density & \cref{ch:maxwell} \\
$\vH$ & magnetic field intensity, $\vH=\nu\vB$ & \cref{ch:maxwell} \\
$\vD$ & electric displacement, $\vD=\varepsilon\vE$ & \cref{ch:maxwell} \\
$\vJ$ & (free) current density & \cref{ch:maxwell} \\
$\vA$ & magnetic vector potential, $\vB=\Curl\vA$ & \cref{ch:magnetostatics} \\
$V$ & electrostatic scalar potential, $\vE=-\Grad V$ & \cref{ch:electrostatics} \\
$\varepsilon$ & electric permittivity & \cref{ch:maxwell} \\
$\mu$ & magnetic permeability & \cref{ch:maxwell} \\
$\nu$ & magnetic reluctivity, $\nu=\mu^{-1}$ & \cref{ch:weak} \\
$\sigma$ & electrical conductivity (loss term) & \cref{ch:maxwell} \\
\bottomrule
\end{tabular}
\end{table}

\begin{table}[ht]
\centering
\small
\caption[Scalar outputs and modal quantities]{Frequencies, modal quantities, and
the scalar products a driver reports.}
\label{tab:gl-outputs}
\begin{tabular}{@{}lll@{}}
\toprule
Symbol & Meaning & Where introduced \\
\midrule
$\omega$ & angular frequency, $\omega=2\pi f$ & \cref{ch:maxwell} \\
$f$ & (ordinary) frequency & \cref{ch:driven} \\
$\lambda$ & eigenvalue (here $\lambda=\omega^2$, or a shifted variant) & \cref{ch:eigenvalues} \\
$k_n$ & wavenumber / propagation constant of mode $n$ & \cref{ch:waveports} \\
$Q$ & quality factor of a resonant mode, $Q=\omega/(2\,\Im\omega)$ & \cref{ch:eigenmodes} \\
$S_{ij}$ & scattering-matrix entry (port $j$ in, port $i$ out) & \cref{ch:driven} \\
$C$ & capacitance (matrix) & \cref{ch:electrostatics} \\
$M$ & inductance (matrix) & \cref{ch:magnetostatics} \\
$U$ & stored field energy, electric or magnetic & \cref{ch:reductions-pde} \\
\bottomrule
\end{tabular}
\end{table}

\begin{notationbox}
The letter $\sigma$ does double duty: it is the conductivity in
\cref{tab:gl-fields} (a physical material property) and the spectral
\emph{shift} of a shift-and-invert eigensolve in \cref{tab:gl-numerics} (a purely
algebraic device). The context---a constitutive law versus an eigenvalue
transform---always disambiguates, and the two never appear in the same equation.
\end{notationbox}

\subsection{Discrete operators and matrices}

Assembling a weak form over a mesh turns each continuous bilinear form into a
sparse matrix. These are set sans-serif to mark them as discrete operators acting
on vectors of degrees of freedom, as opposed to the continuous differential
operators of \cref{tab:gl-spaces}.

\begin{table}[ht]
\centering
\small
\caption[Discrete operators]{The assembled matrices and the maps between
discretizations. Sans-serif marks a discrete (matrix) operator.}
\label{tab:gl-operators}
\begin{tabular}{@{}lll@{}}
\toprule
Symbol & Meaning & Where introduced \\
\midrule
$\mat{K}$ & stiffness operator (curl--curl or grad--grad) & \cref{ch:weak} \\
$\mat{M}$ & mass operator & \cref{ch:weak} \\
$\mat{C}$ & damping / conductivity operator & \cref{ch:weak} \\
$\mat{A}$ & a generic system matrix, $\mat{A}\vx=\vb$ & \cref{ch:cg} \\
$\mat{A}(\omega)$ & frequency-dependent system matrix of a driven solve & \cref{ch:driven} \\
$\mat{P}$ & prolongation (coarse$\to$fine) or a projector & \cref{ch:multigrid} \\
$\mat{G}$ & discrete gradient (the $\Hone\to\Hcurl$ map); a restriction & \cref{ch:smoothers} \\
$a(u,v)$ & a bilinear (weak) form, the source of $\mat{K},\mat{M},\mat{C}$ & \cref{ch:weak} \\
\bottomrule
\end{tabular}
\end{table}

The driven system matrix is the recurring composite
$\mat{A}(\omega)=\mat{K}+i\omega\mat{C}-\omega^2\mat{M}$: one matrix-valued
function of frequency, built once from the three assembled operators and
re-evaluated across a sweep (\cref{ch:driven}). The discrete gradient $\mat{G}$
is the algebraic shadow of $\Grad$---it is what makes the Hiptmair smoother of
\cref{ch:smoothers} possible, by carrying a residual from the edge space into the
nodal space where the near-null-space lives.

\subsection{Function spaces and the finite-element vocabulary}

The four de Rham spaces and the operators that connect them are the backbone of
\cref{part:fieldtheory}; the discretization symbols ($N$, $\eta_K$, $\theta$) and
the element-local axes come from \cref{part:machinery}.

\begin{table}[ht]
\centering
\small
\caption[Function spaces and continuous operators]{The de Rham complex and its
differential operators.}
\label{tab:gl-spaces}
\begin{tabular}{@{}lll@{}}
\toprule
Symbol & Meaning & Where introduced \\
\midrule
$\Hone$ & scalar potentials; nodal elements & \cref{ch:derham} \\
$\Hcurl$ & curl-conforming (edge) elements; tangential continuity & \cref{ch:derham} \\
$\Hdiv$ & div-conforming (face) elements; normal continuity & \cref{ch:derham} \\
$\Ltwo$ & square-integrable fields; no continuity imposed & \cref{ch:derham} \\
$\Grad$ & gradient, $\Hone\to\Hcurl$ & \cref{ch:derham} \\
$\Curl$ & curl, $\Hcurl\to\Hdiv$ & \cref{ch:derham} \\
$\Div$ & divergence, $\Hdiv\to\Ltwo$ & \cref{ch:derham} \\
$\diffop$ & a generic first-order differential operator in a weak form & \cref{ch:weak} \\
\bottomrule
\end{tabular}
\end{table}

\begin{table}[ht]
\centering
\small
\caption[Discretization and adaptivity symbols]{Degree-of-freedom counts, error
indicators, and the element-local axes.}
\label{tab:gl-fem}
\begin{tabular}{@{}lll@{}}
\toprule
Symbol & Meaning & Where introduced \\
\midrule
$N$ & number of (true) degrees of freedom & \cref{ch:fem} \\
$\eta_K$ & local error indicator on element $K$ & \cref{ch:amr} \\
$\theta$ & D\"orfler marking fraction & \cref{ch:amr} \\
E & element axis (which element) & \cref{ch:refelement} \\
L & local-node axis (dof within an element) & \cref{ch:refelement} \\
P & quadrature-point axis & \cref{ch:matrixfree} \\
C & component axis (vector-field component) & \cref{ch:refelement} \\
G & global-dof axis (after gather/scatter) & \cref{ch:assembly} \\
\bottomrule
\end{tabular}
\end{table}

The five element-local axes \texttt{E/L/P/C/G} are the named shape run of the
matrix-free operator: \cref{ch:matrixfree} writes a quadrature contraction as a
sequence of moves along these axes, and \cref{ch:assembly} shows the gather and
scatter that fold the local axis \texttt{L} into the global axis \texttt{G}.

\subsection{Numerical-method symbols}

The Krylov, eigensolver, and multigrid chapters of \cref{part:machinery} share a
compact algebraic vocabulary. The same letters recur---$\alpha$ and $\beta$ for
step coefficients, $\mat{H}$ and $\mat{T}$ for the small projected operators,
$\sigma$ for a spectral shift---so they are collected here.

\begin{table}[ht]
\centering
\small
\caption[Krylov and eigensolver symbols]{Subspaces, coefficients, and projected
operators of the iterative methods.}
\label{tab:gl-numerics}
\begin{tabular}{@{}lll@{}}
\toprule
Symbol & Meaning & Where introduced \\
\midrule
$\mathcal{K}_m$ & Krylov subspace of dimension $m$ & \cref{ch:krylov} \\
$\alpha$ & CG step length / Lanczos diagonal coefficient & \cref{ch:cg} \\
$\beta$ & CG conjugacy factor / Lanczos off-diagonal coefficient & \cref{ch:cg} \\
$\bar{\mat{H}}$ & upper-Hessenberg matrix (Arnoldi / GMRES) & \cref{ch:gmres} \\
$\mat{T}$ & tridiagonal matrix (Lanczos) & \cref{ch:krylovschur} \\
$\sigma$ & spectral shift in shift-and-invert & \cref{ch:eigenvalues} \\
$\kappa$ & condition number, $\kappa(\mat{A})$ & \cref{ch:cg} \\
$c,s$ & cosine / sine of a Givens rotation & \cref{ch:givens} \\
$\mat{V}$ & matrix of (orthonormal) basis vectors & \cref{ch:krylov} \\
\bottomrule
\end{tabular}
\end{table}

The Hessenberg $\bar{\mat{H}}$ and the tridiagonal $\mat{T}$ are the same object
seen from two angles: $\mat{T}$ is what $\bar{\mat{H}}$ collapses to when the
operator is symmetric, which is exactly why Lanczos is the cheap special case of
Arnoldi (\cref{ch:krylovschur}). The Givens pair $(c,s)$ is the rotation that
GMRES applies to keep $\bar{\mat{H}}$ triangular incrementally
(\cref{ch:givens}); the \emph{V-cycle} (\cref{ch:multigrid}) is the recursive
coarsen--smooth--correct schedule that the smoothers of \cref{ch:smoothers}
plug into.

\subsection{The pseudo-language and calculus}

These are the symbols of the typed pseudo-language in which the book writes its
algorithms. They always appear in a shaded box. The combinators and named
operators are the verbs of the calculus; the records are its nouns; the
\texttt{Solve} monad threads the one piece of mutable state.

\begin{table}[ht]
\centering
\small
\caption[Calculus types and shapes]{The type and shape vocabulary of the
pseudo-language.}
\label{tab:gl-calc-types}
\begin{tabular}{@{}lll@{}}
\toprule
Identifier & Meaning & Where introduced \\
\midrule
\texttt{Tensor[N]} & flat vector of $N$ dofs & \cref{ch:tensors} \\
\texttt{Tensor[(S:\,\dots)]} & tensor with a named, rank-agnostic shape run \texttt{S} & \cref{ch:shapes} \\
\texttt{Tensor[\$S]} & a back-reference: ``the same shape \texttt{S} again'' & \cref{ch:shapes} \\
\texttt{Op[\dots]} & a (matrix-free) linear-operator value & \cref{ch:matrixfree} \\
\texttt{Scalar} & a rank-0 tensor (a number with shape) & \cref{ch:tensors} \\
\bottomrule
\end{tabular}
\end{table}

The shape run is the heart of the shape system (\cref{ch:shapes}): a function
typed \texttt{Tensor[(S:\,\dots)] -> Tensor[\$S]} works on a vector, a matrix, or
a higher array at once, and the \texttt{\$S} back-reference promises the output
carries the same axes as the input. We write the sigil form only inside the
shaded box; in prose we name it ``the shape group \texttt{S}.''

\begin{table}[ht]
\centering
\small
\caption[Combinators and named operators]{The loop combinator and the named solve
verbs of the synthesis surface.}
\label{tab:gl-calc-ops}
\begin{tabular}{@{}lll@{}}
\toprule
Identifier & Meaning & Where introduced \\
\midrule
\texttt{iterate\_while} & the one loop combinator (predicate + step) & \cref{ch:fold} \\
\texttt{fold\_solve} & a linear solve expressed as a fold over Krylov steps & \cref{ch:krylov} \\
\texttt{ksp\_solve} & the preconditioned Krylov solve verb & \cref{ch:precond} \\
\texttt{solve\_family} & one operator, many right-hand sides (reused factorization) & \cref{ch:electrostatics} \\
\texttt{frequency\_sweep} & a driven solve evaluated across a band of $\omega$ & \cref{ch:driven} \\
\texttt{eigsolve} & the eigenvalue-problem solve verb & \cref{ch:eigenvalues} \\
\bottomrule
\end{tabular}
\end{table}

\begin{table}[ht]
\centering
\small
\caption[Records and effects]{The state records, the \texttt{Solve} monad, and
the two edge kinds of the dependency graph.}
\label{tab:gl-calc-records}
\begin{tabular}{@{}lll@{}}
\toprule
Identifier & Meaning & Where introduced \\
\midrule
\texttt{IoData} & the parsed configuration record (the run's inputs) & \cref{ch:records} \\
\texttt{OpParams} & the assembled-operator parameter bundle & \cref{ch:records} \\
\texttt{SimState} & the evolving simulation state the monad threads & \cref{ch:monad} \\
\texttt{SolveResult} & a solve's output record (solution + diagnostics) & \cref{ch:records} \\
\texttt{Solve} & the state monad threading \texttt{SimState} & \cref{ch:monad} \\
\texttt{\#extern} & an opaque external kernel (no body; type only) & \cref{ch:matrixfree} \\
\texttt{depends-on} & a blocking edge: constrains resolution and liveness & \cref{ch:wholemachine} \\
\texttt{reference} & a navigational (free) edge: a cross-link, not a dependency & \cref{ch:wholemachine} \\
\bottomrule
\end{tabular}
\end{table}

The \texttt{Solve} monad is the one place state changes: a \texttt{do} block in
the \texttt{Solve} monad threads a \texttt{SimState} through a sequence of steps,
and \texttt{modify} is the single state-mutating verb (\cref{ch:monad}). An
\texttt{\#extern} marks a leaf the synthesis surface does not implement---a call
out to an opaque library kernel (libCEED quadrature, a vendor eigensolve), shown
by its type signature only (\cref{ch:matrixfree}).

\subsection{The layered stack and methodology}

The book's organizing idea is a stack of representations connected by rotations.
These tokens name its levels and its recurring structural patterns.

\begin{table}[ht]
\centering
\small
\caption[The layered stack]{Levels, rotations, and the structural categories of
the methodology.}
\label{tab:gl-method}
\begin{tabular}{@{}lll@{}}
\toprule
Token & Meaning & Where introduced \\
\midrule
L0 & ground-truth source level (cited ranges) & \cref{ch:strata} \\
L1 & pure-function level (mutation rotated out) & \cref{ch:strata} \\
L2 & fused-algebra level (HPC tricks erased) & \cref{ch:strata} \\
L3 & iteration level (global tensor-field ops) & \cref{ch:strata} \\
L4 & calculus level (combinators + monad + tensors) & \cref{ch:strata} \\
rotation & a representation-preserving rewrite between levels & \cref{ch:rotations} \\
obstruction & a place a rotation cannot complete cleanly & \cref{ch:rotations} \\
\bottomrule
\end{tabular}
\end{table}

A \emph{rotation} re-expresses a level in a new vocabulary while preserving what
it computes; an \emph{obstruction} is an honest record of a place that resists
the rotation (a sequential recurrence that will not become a tensor-field op, an
opaque library boundary). Two small taxonomies recur. The \emph{four solve
corners} (\cref{ch:composing}) are the cross product of two axes---one
right-hand side versus a family, and a fixed operator versus a frequency-varying
one---and they classify the five drivers: electrostatic and magnetostatic are
family-of-RHS / fixed-operator corners, the driven solve is the
varying-operator corner, and so on. The \emph{three reduction shapes}
(\cref{ch:reductions-pde}) are the three ways a field is collapsed to a
reported number---a global scalar (energy), a boundary integral (a port flux),
and a pointwise probe (a field sample at a point)---one shape per family of
output product.

\begin{keyidea}
Every symbol in this glossary lives at a definite level of the stack. The
physical fields ($\vE$, $\vB$) and the function spaces ($\Hcurl$, $\Ltwo$) name
the continuous problem; the assembled operators ($\mat{K}$, $\mat{M}$) and the
Krylov symbols ($\alpha$, $\bar{\mat{H}}$) name the discrete machinery; the
pseudo-language identifiers (\texttt{iterate\_while}, \texttt{Solve}) name the
calculus that rewrites that machinery into a synthesizable surface. Reading a
symbol's level is half of reading the symbol.
\end{keyidea}

\section{A guide to the literature}
\label{sec:furtherreading-app}

The book is a synthesis, and it stands on a small, deliberately chosen
bibliography. This section groups those sources by topic and says, in a line,
what to read each for. It is a map for going deeper, not a survey.

\paragraph{Finite elements and Maxwell's equations.}
For the finite-element treatment of electromagnetics, \citet{monk2003} is the
standard reference and the closest in spirit to the field-theory part of this
book---read it for the rigorous theory of edge elements and $\Hcurl$
conformity. \citet{jin2014} is the more applied companion, rich in engineering
detail and worked formulations; read it for how the weak forms become real
solvers. \citet{brennerscott2008} is the mathematical foundation of the finite
element method in general (approximation theory, error estimates); read it when
you want the analysis under the discretization. \citet{arnold2006feec} is the
deep structural account: finite element exterior calculus explains \emph{why}
the de Rham complex of \cref{ch:derham} is the right organizing principle, and
why the four spaces must connect as they do.

\paragraph{Krylov methods and iterative linear solvers.}
\citet{saad2003} is the encyclopedic reference for sparse iterative
methods---read it for CG, GMRES, preconditioning, and the Krylov-subspace theory
behind \cref{ch:krylov,ch:cg,ch:gmres}. \citet{trefethenbau1997} is the elegant
short course; read it for the cleanest possible derivations of QR, Arnoldi, and
the conditioning ideas ($\kappa$). \citet{golubvanloan2013} is the matrix-computations
handbook; read it for Givens rotations (\cref{ch:givens}), the dense linear
algebra under the small projected problems, and as a definitive cross-reference.

\paragraph{Eigensolvers.}
\citet{saad2011eig} is the companion volume to Saad's linear-systems book and the
reference for large-scale eigenvalue computation---read it for Lanczos,
Arnoldi, and the shift-and-invert ($\sigma$) transforms of
\cref{ch:eigenvalues}. \citet{stewart2001krylovschur} introduced the
Krylov--Schur restarting scheme of \cref{ch:krylovschur}; read it for the
stable thick-restart mechanism. \citet{effenberger2013} treats the nonlinear
eigenvalue problem and the successive-deflation strategy behind
\cref{ch:deflation}; read it when the eigenproblem's operator depends on its own
eigenvalue.

\paragraph{Multigrid and preconditioning.}
\citet{hiptmair1998} is the foundational paper on multigrid for Maxwell's
equations and the source of the Hiptmair smoother (\cref{ch:smoothers})---read
it for why curl-conforming problems need a special near-null-space correction.
\citet{hiptmairxu2007} generalizes this into nodal auxiliary-space
preconditioning (the AMS idea); read it for the auxiliary-space construction
that makes $\Hcurl$ and $\Hdiv$ solves scalable. \citet{phillipsfischer2022}
gives the optimal Chebyshev smoothers and one-sided V-cycle schedule of
\cref{ch:multigrid}; read it for the practical smoother tuning.

\paragraph{Adaptivity and error estimation.}
\citet{dorfler1996} introduced the marking strategy ($\theta$) that drives the
adaptive loop of \cref{ch:amr}---read it for the convergence guarantee behind
``refine the smallest set of elements holding a fixed fraction of the error.''
\citet{zienkiewiczzhu1987} is the classic recovery-based error estimator
($\eta_K$); read it for the practical, inexpensive indicator that tells the loop
\emph{where} to refine.

\paragraph{Microwave engineering and physics.}
\citet{pozar2011} is the microwave-engineering reference and the source of the
S-parameter and wave-port conventions ($S_{ij}$, $k_n$) of
\cref{ch:driven,ch:waveports}---read it for what the driven solve's outputs
mean physically. \citet{griffiths2017} is the undergraduate electrodynamics
text; read it to refresh Maxwell's equations, the potentials, and the energy
integrals of \cref{ch:maxwell} from first principles.

\paragraph{Functional programming and the calculus.}
\citet{peytonjones2003} is the Haskell language report and the reference for the
typed-functional notation the pseudo-language borrows (currying, records, type
signatures)---read it for the precise meaning of the syntax in the shaded
boxes. \citet{wadler1995monads} is the classic exposition of monads for
functional programming; read it for the idea that makes the \texttt{Solve} monad
of \cref{ch:monad} legitimate---state threaded purely, with one mutation point.

\paragraph{The tools.}
\citet{palace2023} is Palace itself, the AWS Labs finite-element electromagnetics
solver this book dissects; it is the L0 ground truth that every citation in the
book ultimately points at. \citet{ceed2021} is libCEED, the matrix-free
high-order element library behind the \texttt{\#extern} quadrature kernel of
\cref{ch:matrixfree}; read it for the element-local algebra (the E/L/P/C/G axes)
the synthesis surface lifts into tensor contractions.

\summarysection
You now have the whole book in three forms. You have the \emph{map}---the
layered stack L0 through L4, the rotations that connect the levels, and the
honest obstructions where a rotation cannot finish. You have the
\emph{machinery}---the weak forms and de Rham spaces of the field theory, the
Krylov, eigensolver, and multigrid algorithms of the numerical core, and the
five drivers that compose them into electrostatic, magnetostatic, eigenmode,
driven, and transient simulations. And you have the \emph{vocabulary}---the
symbols of this glossary and the typed pseudo-language they belong to, in which
the synthesis surface re-expresses a quarter-million lines of tuned C++ as a few
combinators over immutable tensors. The simulator is no longer a black box: it is
a small calculus, a handful of named operators, and a clear account of what each
one computes and why. That account was the goal; you are now holding it.
